\chapter{Building and running the program}
\label{chap-utilisation}

\section{Prerequisites}

A C11 compiler, CMake~3.25 or later, \code{pkg-config} and the development files of GTK~3.
GStreamer is optional (sound effects).

\begin{table}[ht]
    \centering
    \small
    \begin{tabularx}{\textwidth}{@{}l>{\raggedright\arraybackslash}X@{}}
        \toprule
        \textbf{System}  & \textbf{Installing the dependencies}                                \\
        \midrule
        Debian, Ubuntu   & \code{sudo apt install build-essential cmake pkg-config libgtk-3-dev
                               libgstreamer1.0-dev} \\
        Fedora           & \code{sudo dnf install gcc cmake pkgconf-pkg-config gtk3-devel
                               gstreamer1-devel}       \\
        macOS            & \code{brew install cmake pkg-config gtk+3 gstreamer}                \\
        \bottomrule
    \end{tabularx}
\end{table}

\section{Commands}

\begin{lstlisting}[language=bash, morekeywords={make, cmake, ctest, git}]
git clone https://github.com/DmytroPalahin/battleship.git
cd battleship

make test        # configure, compile, run the tests
make run         # start the game
make play        # play in the terminal
make simulate    # measure the three levels (10,000 games)
make asan        # tests under AddressSanitizer and UBSan
make report      # build this report
make help        # complete list
\end{lstlisting}

The \fichier{Makefile} is only a reminder: it calls the CMake presets. The explicit form is
\code{cmake -{}-preset release}, \code{cmake -{}-build -{}-preset release},
\code{ctest -{}-preset release}.

\section{Command-line options}

\begin{table}[ht]
    \centering
    \small
    \begin{tabularx}{\textwidth}{@{}l>{\raggedright\arraybackslash}X@{}}
        \toprule
        \multicolumn{2}{@{}l}{\textbf{\code{battleship}} --- graphical application}               \\
        \midrule
        \code{-l}, \code{-{}-level=LEVEL}  & Start directly at a level: \code{easy},
        \code{medium} or \code{hard}.                                                             \\
        \code{-s}, \code{-{}-seed=N}       & Seed of the first game, to replay it.                \\
        \code{-m}, \code{-{}-mute}         & Turn the sound effects off.                          \\
        \code{-V}, \code{-{}-version}      & Print the version.                                   \\
        \code{-{}-autoplay}                & Test: the computer plays both sides, then the
        program closes.                                                                           \\
        \code{-{}-max-shots=N}             & Test: with \code{-{}-autoplay}, stop after $N$
        shots.                                                                                    \\
        \code{-{}-screenshot=FILE}         & Test: save the window as a PNG file before
        quitting.                                                                                 \\
        \addlinespace
        \multicolumn{2}{@{}l}{\textbf{\code{battleship-cli}} --- terminal interface}              \\
        \midrule
        \code{play}                        & Play against the computer (\code{-{}-level},
        \code{-{}-seed}).                                                                         \\
        \code{simulate}                    & Measure the levels (\code{-{}-games},
        \code{-{}-seed}, \code{-{}-level}, \code{-{}-csv},
        \code{-{}-no-duel}).                                                                      \\
        \bottomrule
    \end{tabularx}
\end{table}

\section{Controls}

With the mouse: click on a cell of the pirate fleet. With the keyboard: the arrow keys move
the missile over the enemy grid, \touche{Space} fires, \touche{F11} toggles full screen.
